% ===================================================================
% Section 6: Applications
% ===================================================================
\section{Application to US sector ETFs}
\label{sec:applications}

\subsection{Data and protocol}
We evaluate on a real universe of $11$ US sector exchange-traded funds (XLB,
XLC, XLE, XLF, XLI, XLK, XLP, XLRE, XLU, XLV, XLY), using daily total returns
from July 2018 onwards, roughly $2000$ observations per asset. The backtest is a
rolling out-of-sample exercise: at each rebalancing date the covariance is
estimated on a trailing window of $252$ trading days (one year), the portfolio is
solved on that estimate, and the resulting weights are held, unchanged and
out-of-sample, until the next rebalance $21$ trading days (approximately one
month) later. This yields $84$ rebalances. No information from the holding period
enters the estimate.

We compare six strategies: equal weight; sample-covariance Markowitz (GMV);
a constrained GMV with a per-asset cap; ridge-GMV; Ledoit--Wolf shrinkage GMV;
and the proposed design-regularized portfolio, which uses
$\lambda_2=0.10$, $\gD=\gE=0.05$, long-only with a $30\%$ cap, and the
precision-metric specification $V=\Sighat^{-1/2}$ of
Example~\ref{ex:Vchoices}(2). Metrics are the annualized Sharpe ratio,
out-of-sample annualized volatility, average turnover per rebalance, average
effective number of assets $\Neff$, and maximum drawdown.

\subsection{Results}

\begin{table}[htbp]
\centering
\caption{Real sector-ETF backtest, July 2018 onwards; $252$-day estimation
  window, monthly rebalancing, $84$ rebalances. Lower turnover and higher
  effective $N$ are better. ``Proposed'' is ridge $+$ D $+$ E with a $30\%$
  per-asset cap.}
\label{tab:real}
\begin{tabular}{lrrrrr}
\toprule
\textbf{Strategy} & \textbf{Sharpe} & \textbf{OOS vol} & \textbf{Turnover} & \textbf{Eff.\ $N$} & \textbf{Max DD} \\
\midrule
Equal weight             & \textbf{0.751} & 0.1875 & 0.000 & 11.00 & $-0.363$ \\
Markowitz (sample)       & 0.483 & 0.1573 & 0.423 & 1.68 & $-0.338$ \\
Constrained GMV          & 0.651 & 0.1632 & 0.116 & 4.18 & $-0.332$ \\
Ridge GMV                & 0.655 & 0.1734 & 0.063 & 8.67 & $-0.357$ \\
Shrinkage GMV            & 0.528 & 0.1555 & 0.311 & 2.19 & $-0.343$ \\
Proposed (D$+$E$+$ridge) & 0.710 & 0.1766 & \textbf{0.038} & \textbf{9.30} & $-0.355$ \\
\bottomrule
\end{tabular}
\end{table}

Reading Table~\ref{tab:real} honestly:

\begin{itemize}[leftmargin=1.4em,itemsep=0.25em]
  \item The raw sample Markowitz portfolio is the worst risk-adjusted performer
    among the optimized strategies (Sharpe $0.483$), concentrates into fewer than
    two effective names out of eleven, and churns heavily (turnover $0.423$).
    This is precisely the instability that motivates the paper, and it persists
    even in a small, well-conditioned, liquid universe.
  \item The \textbf{proposed} portfolio recovers almost all of the equal-weight
    risk-adjusted return ($0.710$ against $0.751$) while cutting turnover by
    roughly an order of magnitude relative to sample Markowitz ($0.038$ against
    $0.423$) and raising effective breadth from $1.68$ to $9.30$ names. Among the
    optimized strategies it dominates on the stability-for-return trade-off.
  \item Ridge alone already helps substantially (turnover $0.063$, $\Neff=8.67$);
    the design penalties push turnover lower still and breadth higher, at a small
    cost in out-of-sample volatility. This is the same complementarity seen in
    the high-dimensional simulation of Section~\ref{sec:highdim}.
  \item Shrinkage improves the covariance estimate---it attains the lowest
    out-of-sample volatility---but, because it does not regularize the
    \emph{weights}, it does not by itself control turnover or concentration.
    This separation between estimator regularization and weight regularization
    is the practical case for the present formulation.
\end{itemize}

The turnover ordering across strategies is consistent with
Corollary~\ref{cor:turnover}: the strategies whose objectives carry a larger
strong-convexity modulus---ridge, and then ridge plus barrier penalties---exhibit
systematically lower turnover, while the two strategies with no weight-space
regularization (sample Markowitz and shrinkage) sit at the top of the turnover
ranking.

\begin{figure}[htbp]
  \centering
  \includegraphics[width=\textwidth]{week8_metrics.png}
  \caption{Real sector-ETF metric panel: annualized Sharpe ratio, out-of-sample
    volatility, average turnover, and average effective $N$ by strategy. The
    proposed portfolio combines equal-weight-like breadth with the lowest
    turnover among the optimized methods.}
  \label{fig:realmetrics}
\end{figure}

\begin{figure}[htbp]
  \centering
  \includegraphics[width=0.85\textwidth]{week8_equity_curves.png}
  \caption{Out-of-sample cumulative growth of one dollar on the sector-ETF
    universe, July 2018 onwards.}
  \label{fig:realeq}
\end{figure}

\subsection{Transaction costs}
Because turnover translates directly into trading cost, the turnover reduction
is economically and not merely statistically meaningful. Under a linear cost
model charging a fixed number of basis points per unit of turnover, the
net-of-cost Sharpe ranking increasingly favours the low-turnover
design-regularized portfolio as the cost level rises, whereas the raw Markowitz
portfolio's heavy churn erodes its already weak gross performance. At zero cost
the ranking is that of Table~\ref{tab:real}; as costs increase, the gap widens
monotonically in favour of the proposed strategy, which is the practical payoff
of Corollary~\ref{cor:turnover}.

\subsection{Crisis and calm sub-periods}
Splitting the sample into stressed and calm sub-periods gives a mixed but
interpretable read on the E-criterion. In the calm regime the robust, E-tilted
variant improves risk-adjusted return relative to GMV, while in the stressed
regime its realized volatility is comparable to that of GMV and its maximum
drawdown is essentially unchanged. We therefore do not claim a crisis-robustness
result. The worst-case payoff that E-optimality is designed to deliver should be
clearest under sharper and more genuinely out-of-sample stress than a
sub-period split of a single recent sample can provide, and we flag this as the
main direction for further empirical work rather than as an established finding.
